\section{Bài tập 10}

Cho mạch điện như hình bên dưới. Hãy tìm điện áp $V$. Có thể giải bằng bất kỳ phương pháp nào nhưng cần giải thích chi tiết. Sau đó, mô phỏng mạch điện để kiểm tra kết quả.

\given{ex10_given}{0.55\textwidth}{Tìm giá trị điện áp $V$}

\subsection{Tính toán}
\noindent\textbf{Phương pháp: Biến đổi Tam giác -- Sao ($\Delta \to Y$) và Quy tắc chia áp}
\begin{enumerate}[leftmargin=2em]
  \item \textbf{Biến đổi $\Delta \to Y$ tam giác $B$-$C$-$D$:} \\
        Tam giác gồm ba điện trở $R_6 = 15\,\Omega$ ($B$-$C$), $R_5 = 10\,\Omega$ ($C$-$D$), và $R_7 = 30\,\Omega$ ($B$-$D$).
        \[
        \Sigma_\Delta = R_6 + R_5 + R_7 = 15 + 10 + 30 = 55\,\Omega
        \]
        Các điện trở hình sao tương đương nối về tâm sao $N$:
        \begin{align*}
        R_B &= \frac{R_6 \cdot R_7}{\Sigma_\Delta} = \frac{15 \times 30}{55} = \frac{90}{11}\,\Omega \approx 8{,}182\,\Omega \\
        R_C &= \frac{R_6 \cdot R_5}{\Sigma_\Delta} = \frac{15 \times 10}{55} = \frac{30}{11}\,\Omega \approx 2{,}727\,\Omega \\
        R_D &= \frac{R_5 \cdot R_7}{\Sigma_\Delta} = \frac{10 \times 30}{55} = \frac{60}{11}\,\Omega \approx 5{,}455\,\Omega
        \end{align*}

  \item \textbf{Gộp hai nhánh song song từ $N$ xuống đất $E$ ($V_E = 0\text{ V}$):}
        \begin{align*}
        R_{NE1} &= R_C + R_3 = 2{,}727\,\Omega + 12\,\Omega = 14{,}727\,\Omega \\
        R_{NE2} &= R_D + R_4 = 5{,}455\,\Omega + 20\,\Omega = 25{,}455\,\Omega \\
        R_{NE} &= R_{NE1} \parallel R_{NE2} = \frac{14{,}727 \times 25{,}455}{14{,}727 + 25{,}455} \approx 9{,}330\,\Omega
        \end{align*}

  \item \textbf{Tổng trở nhánh qua $N$ và điện trở tương đương $R_{BE}$:}
        \begin{align*}
        R_{BNE} &= R_B + R_{NE} = 8{,}182\,\Omega + 9{,}330\,\Omega = 17{,}511\,\Omega \\
        R_{BE} &= R_2 \parallel R_{BNE} = \frac{35 \times 17{,}511}{35 + 17{,}511} \approx 11{,}672\,\Omega
        \end{align*}

  \item \textbf{Tính điện áp $V$ bằng quy tắc chia áp giữa $R_{BE}$ và $R_1$:}
        \[
        V = 100\text{ V} \times \frac{R_{BE}}{R_1 + R_{BE}} = 100\text{ V} \times \frac{11{,}672\,\Omega}{16\,\Omega + 11{,}672\,\Omega} = 100 \times \frac{11{,}672}{27{,}672} \approx 42{,}18\text{ V}
        \]
\end{enumerate}

\begin{dapsobox}
  \begin{itemize}[leftmargin=1.5em]
    \item $V \approx 42{,}18\text{ V}$
  \end{itemize}
\end{dapsobox}

\subsection{Mô phỏng}
\simresult
\simfig{ex10_sim}{Kết quả mô phỏng PSpice của Bài tập 10}
